\documentclass[11pt,a4paper]{article}

% ---------------------------------------------------------
% PREAMBLE (stable, warning-free, Overleaf-safe)
% ---------------------------------------------------------
\usepackage[utf8]{inputenc}
\usepackage[T1]{fontenc}
\usepackage{lmodern}
\usepackage{amsmath, amssymb, amsfonts}
\usepackage{bm}
\usepackage{geometry}
\usepackage{graphicx}
\usepackage{hyperref}
\usepackage{cite}
\usepackage{authblk}
\usepackage{physics}
\usepackage{mathtools}

\geometry{margin=2.4cm}

\title{\textbf{Companion LI: The Time Derivative of the Gravitational Potential in the Relaxation-Driven Cyclic Cosmology}}
\author{Michael Lehmann \\ \texttt{mi.lehmann@gmx.de}}
\date{April 2026}

\begin{document}
\maketitle

\begin{abstract}
The time derivative of the gravitational potential, $\dot{\Phi}$, plays a central 
role in several cosmological observables, including the integrated Sachs--Wolfe 
(ISW) effect, the Rees--Sciama effect, and the evolution of large-scale structure. 
In the standard cosmological model, $\dot{\Phi}$ is small during matter domination 
and becomes significant only at late times due to dark energy. In the 
Relaxation-Driven Cyclic Cosmology (RDCC), the scale-dependent evolution of the 
gravitational potential introduces a new temporal structure in $\dot{\Phi}$, as the 
relaxation mechanism modifies the decay of small-scale modes and enhances the 
coherence of large-scale modes. This Companion develops a continuous description of 
$\dot{\Phi}$ in RDCC, deriving how the relaxation scale influences the ISW signal, 
the Rees--Sciama effect, and the temporal evolution of the gravitational field. The 
resulting framework provides a direct observational probe of the RDCC relaxation 
mechanism.
\end{abstract}
% ---------------------------------------------------------
\section{Introduction}

The gravitational potential $\Phi$ governs the motion of matter and the propagation 
of light. Its time derivative, $\dot{\Phi}$, encodes the dynamical evolution of the 
gravitational field and plays a central role in several cosmological observables. 
The integrated Sachs--Wolfe (ISW) effect arises from the line-of-sight integral of 
$\dot{\Phi}$, while the Rees--Sciama effect arises from non-linear contributions to 
$\dot{\Phi}$.

In the standard cosmological model, $\dot{\Phi}$ is small during matter domination 
and becomes significant only at late times due to dark energy. In RDCC, the 
gravitational potential evolves differently. The relaxation mechanism suppresses the 
decay of small-scale modes and enhances the coherence of large-scale modes. This 
modifies the temporal structure of $\dot{\Phi}$ in a scale-dependent manner, 
producing distinctive signatures in the ISW and Rees--Sciama effects.
% ---------------------------------------------------------
\section{The RDCC Gravitational Potential}

The gravitational potential in RDCC is given by
\begin{equation}
    \Phi_{\mathrm{RDCC}}(k,a)
    = \Phi(k,a)
      \left[
        1 - \chi_{\Phi}
        \left(\frac{k}{k_{\mathrm{rel}}}\right)^{\alpha}
      \right].
\end{equation}

The time derivative becomes
\begin{equation}
    \dot{\Phi}_{\mathrm{RDCC}}(k,a)
    = \dot{\Phi}(k,a)
      \left[
        1 - \chi_{\Phi}
        \left(\frac{k}{k_{\mathrm{rel}}}\right)^{\alpha}
      \right]
      - \Phi(k,a)\,
        \chi_{\Phi}\,\alpha\,
        \left(\frac{k}{k_{\mathrm{rel}}}\right)^{\alpha}
        \frac{\dot{k}_{\mathrm{rel}}}{k_{\mathrm{rel}}}.
\end{equation}

This expression shows that $\dot{\Phi}$ acquires a new contribution from the 
relaxation mechanism.
% ---------------------------------------------------------
\section{The RDCC ISW Effect}

The ISW temperature shift is given by
\begin{equation}
    \frac{\Delta T_{\mathrm{ISW}}}{T_{\gamma}}
    = 2 \int \dot{\Phi}\, d\chi.
\end{equation}

In RDCC, this becomes
\begin{equation}
    \Delta T_{\mathrm{ISW,RDCC}}
    = \Delta T_{\mathrm{ISW}}
      \left[
        1 - \chi_{\mathrm{ISW}}
        \left(\frac{k}{k_{\mathrm{rel}}}\right)^{\alpha}
      \right]
      + 2 \int d\chi\,
        \Phi(k,a)\,
        \chi_{\Phi}\,\alpha\,
        \left(\frac{k}{k_{\mathrm{rel}}}\right)^{\alpha}
        \frac{\dot{k}_{\mathrm{rel}}}{k_{\mathrm{rel}}}.
\end{equation}

The second term is unique to RDCC and provides a direct observational probe of the 
relaxation scale.
% ---------------------------------------------------------
\section{The Rees--Sciama Effect in RDCC}

The Rees--Sciama effect arises from non-linear contributions to $\dot{\Phi}$. In 
RDCC, the non-linear potential becomes
\begin{equation}
    \Phi_{\mathrm{NL,RDCC}}
    = \Phi_{\mathrm{NL}}
      \left[
        1 - \chi_{\mathrm{NL}}
        \left(\frac{k}{k_{\mathrm{rel}}}\right)^{\alpha}
      \right].
\end{equation}

The resulting Rees--Sciama temperature shift becomes
\begin{equation}
    \Delta T_{\mathrm{RS,RDCC}}
    = \Delta T_{\mathrm{RS}}
      \left[
        1 - \chi_{\mathrm{RS}}
        \left(\frac{k}{k_{\mathrm{rel}}}\right)^{\alpha}
      \right].
\end{equation}

This modification suppresses small-scale contributions and enhances large-scale 
coherence.
% ---------------------------------------------------------
\section{Temporal Coherence and RDCC Signatures}

The relaxation mechanism modifies the temporal coherence of the gravitational 
potential. The correlation function of $\dot{\Phi}$ becomes
\begin{equation}
    \xi_{\dot{\Phi},\mathrm{RDCC}}(k,a)
    = \xi_{\dot{\Phi}}(k,a)
      \left[
        1 - \chi_{\xi}
        \left(\frac{k}{k_{\mathrm{rel}}}\right)^{\alpha}
      \right].
\end{equation}

This modification produces:

\begin{itemize}
    \item enhanced large-scale ISW correlations,
    \item suppressed small-scale Rees--Sciama fluctuations,
    \item a characteristic turnover near the relaxation scale.
\end{itemize}

These signatures provide a direct observational probe of the RDCC gravitational 
field.
% ---------------------------------------------------------
\section{Conclusion}

The time derivative of the gravitational potential in the Relaxation-Driven Cyclic 
Cosmology reflects the scale-dependent evolution of the gravitational field. The 
relaxation mechanism modifies $\dot{\Phi}$ in a scale-dependent manner, producing 
distinctive signatures in the ISW effect, the Rees--Sciama effect, and the temporal 
coherence of the gravitational potential. These effects provide a direct 
observational window into the relaxation scale and the temporal structure of the 
RDCC gravitational field. The present Companion establishes the theoretical 
foundation for future studies of CMB secondary anisotropies, potential evolution, 
and the interplay between matter and gravity in RDCC.

\bibliographystyle{apsrev4-2}
\nocite{*}
\bibliography{references}

\end{document}
